\chapter{Introducción}

\section{Contexto y Relevancia}
\subsection{Gastronomía como motor cultural y económico}
Las Pequeñas y Medianas Empresas (PYMES) del sector gastronómico desempeñan un papel fundamental en la economía local y en la preservación de la identidad cultural de la ciudad de Portoviejo. La gastronomía no solo constituye una actividad económica relevante, sino también un elemento representativo del patrimonio cultural de la región, reconocido internacionalmente tras la declaratoria de Portoviejo como Ciudad Creativa de la Gastronomía por la UNESCO en 2019. Este reconocimiento resalta el potencial del sector gastronómico como motor de desarrollo económico, generación de empleo y fortalecimiento del emprendimiento local.

\subsection{Estructura empresarial en Portoviejo}
A nivel nacional, las PYMES conforman la mayor parte del tejido empresarial ecuatoriano. Según datos del Observatorio de la PyME, más del 90 \% de las empresas registradas en el país corresponden a microempresas, seguidas por pequeñas y medianas empresas. En la provincia de Manabí, y particularmente en Portoviejo, estas organizaciones cumplen un rol clave en el desarrollo económico y social; sin embargo, muchas de ellas enfrentan limitaciones estructurales que afectan su competitividad y sostenibilidad en un entorno cada vez más digitalizado.

\section{Problemática y Desafíos}
\subsection{Baja adopción de tecnologías digitales}
Uno de los principales desafíos que enfrentan las PYMES gastronómicas es el bajo nivel de adopción de tecnologías digitales. A pesar del creciente reconocimiento de la importancia de la transformación digital, numerosos negocios continúan gestionando sus procesos operativos mediante métodos manuales o herramientas básicas. La administración de pedidos, el control de inventarios, la gestión de repartidores y la atención al cliente suelen realizarse sin sistemas integrados, lo que genera errores operativos, desperdicio de insumos, tiempos de respuesta elevados y dificultades para la toma de decisiones estratégicas.

\subsection{Factores agravantes}
Esta problemática se ve agravada por diversos factores, entre ellos el alto costo de las soluciones tecnológicas disponibles en el mercado, las cuales generalmente están orientadas a empresas medianas o grandes y no se ajustan a la realidad económica de las PYMES locales. Asimismo, la falta de capacitación técnica del personal y la dependencia de plataformas de delivery de terceros, que cobran comisiones elevadas, limitan la rentabilidad y reducen el control que los negocios tienen sobre sus propios procesos. Aunque el Estado ecuatoriano ha impulsado iniciativas como la Agenda de Transformación Digital del Ecuador 2022--2025, que busca fomentar la adopción tecnológica en distintos sectores productivos, la implementación efectiva de estas estrategias en las PYMES gastronómicas aún es limitada.

\section{Brecha y Necesidades}
En este contexto, se identifica una brecha significativa entre el potencial gastronómico de Portoviejo y el nivel real de digitalización de sus pequeñas y medianas empresas. La ausencia de soluciones tecnológicas accesibles, integradas y adaptadas al contexto local impide que estos negocios optimicen sus operaciones, mejoren la experiencia del cliente y compitan en condiciones más equitativas frente a empresas de mayor escala. Esta situación evidencia la necesidad de desarrollar alternativas tecnológicas que respondan de manera específica a las necesidades operativas, técnicas y económicas del sector gastronómico local.

\section{Propuesta de Solución}
Ante esta realidad, el presente estudio propone el desarrollo de una plataforma SaaS (Software como Servicio) integral orientada a las PYMES gastronómicas de la ciudad de Portoviejo. Dicha plataforma busca integrar módulos de gestión de pedidos, control de inventarios, administración de delivery, analítica de datos y atención al cliente automatizada mediante inteligencia artificial, permitiendo a los negocios mejorar su eficiencia operativa sin requerir grandes inversiones en infraestructura tecnológica. El propósito de esta investigación es contribuir a la transformación digital del sector gastronómico local, fortaleciendo la competitividad, sostenibilidad y calidad del servicio ofrecido por las PYMES, en concordancia con las políticas nacionales de digitalización y con la visión de Portoviejo como una ciudad innovadora y creativa \parencite[]{Castells2010}.
